\documentclass[10pt,a4paper]{amsart}
\usepackage{amsmath,amssymb,amsthm,graphicx}
% Retained preamble for the Stacks Project chapter snapshot in this directory.
% The chapter's own preamble.tex is a copy of the shared Stacks Project preamble; the block
% below is sliced MECHANICALLY, line for line, from it. Nothing was added, removed or
% redefined; the wrapper class, the xr-hyper/hyperref wiring and the externaldocument list
% are omitted because the assembler supplies the document class.
%
% lines taken from preamble.tex: 3, 5, 9, 11, 12, 156, 157, 158, ... (51 lines total)
%
% package coverage restored from the chapter's own preamble (2026-10-13): xy, xy-2cell

\usepackage[all]{xy}
\xyoption{2cell}
\UseAllTwocells

\newenvironment{reference}{\comment}{\endcomment}
\newenvironment{slogan}{\comment}{\endcomment}
\newenvironment{history}{\comment}{\endcomment}
\theoremstyle{plain}
\newtheorem{theorem}[subsection]{Theorem}
\newtheorem{proposition}[subsection]{Proposition}
\newtheorem{lemma}[subsection]{Lemma}
\theoremstyle{definition}
\newtheorem{definition}[subsection]{Definition}
\newtheorem{example}[subsection]{Example}
\newtheorem{exercise}[subsection]{Exercise}
\newtheorem{situation}[subsection]{Situation}
\theoremstyle{remark}
\newtheorem{remark}[subsection]{Remark}
\newtheorem{remarks}[subsection]{Remarks}
\numberwithin{equation}{subsection}
\def\lim{\mathop{\mathrm{lim}}\nolimits}
\def\colim{\mathop{\mathrm{colim}}\nolimits}
\def\Spec{\mathop{\mathrm{Spec}}}
\def\Hom{\mathop{\mathrm{Hom}}\nolimits}
\def\Ext{\mathop{\mathrm{Ext}}\nolimits}
\def\SheafHom{\mathop{\mathcal{H}\!\mathit{om}}\nolimits}
\def\SheafExt{\mathop{\mathcal{E}\!\mathit{xt}}\nolimits}
\def\Sch{\mathit{Sch}}
\def\Mor{\mathop{\mathrm{Mor}}\nolimits}
\def\Ob{\mathop{\mathrm{Ob}}\nolimits}
\def\Sh{\mathop{\mathit{Sh}}\nolimits}
\def\NL{\mathop{N\!L}\nolimits}
\def\CH{\mathop{\mathrm{CH}}\nolimits}
\def\proetale{{pro\text{-}\acute{e}tale}}
\def\etale{{\acute{e}tale}}
\def\QCoh{\mathit{QCoh}}
\def\Ker{\mathop{\mathrm{Ker}}}
\def\Im{\mathop{\mathrm{Im}}}
\def\Coker{\mathop{\mathrm{Coker}}}
\def\Coim{\mathop{\mathrm{Coim}}}
\DeclareMathSymbol{\boxtimes}{\mathbin}{AMSa}{"02}
\def\QCohstack{\mathcal{QC}\!\mathit{oh}}
\def\Cohstack{\mathcal{C}\!\mathit{oh}}
\def\Spacesstack{\mathcal{S}\!\mathit{paces}}
\def\Quotfunctor{\mathrm{Quot}}
\def\Hilbfunctor{\mathrm{Hilb}}
\def\Curvesstack{\mathcal{C}\!\mathit{urves}}
\def\Polarizedstack{\mathcal{P}\!\mathit{olarized}}
\def\Complexesstack{\mathcal{C}\!\mathit{omplexes}}
\def\Pic{\mathop{\mathrm{Pic}}\nolimits}
\def\Picardstack{\mathcal{P}\!\mathit{ic}}
\def\Picardfunctor{\mathrm{Pic}}
\def\Deformationcategory{\mathcal{D}\!\mathit{ef}}
\usepackage{cleveref}
\crefname{originallemma}{original lemma}{original lemmas}
\Crefname{originallemma}{Original lemma}{Original lemmas}
\makeatletter
\newlabel{equation-beauville-laszlo-glueing-spaces}{{0F9N}{source}{}{sourcepointer-1}{}}
\newlabel{equation-beauville-laszlo-glueing-spaces@cref}{{[originallemma][0][]0F9N}{[1][0][]source}{}{}{}}
\makeatother
\crefname{originallemma}{original lemma}{original lemmas}
\Crefname{originallemma}{Original lemma}{Original lemmas}
\makeatletter
\newlabel{lemma-glueing-affines}{{0F9Q}{source}{}{sourcepointer-2}{}}
\newlabel{lemma-glueing-affines@cref}{{[originallemma][0][]0F9Q}{[1][0][]source}{}{}{}}
\makeatother
\crefname{originallemma}{original lemma}{original lemmas}
\Crefname{originallemma}{Original lemma}{Original lemmas}
\makeatletter
\newlabel{lemma-glueing-affines-etale}{{0F9R}{source}{}{sourcepointer-3}{}}
\newlabel{lemma-glueing-affines-etale@cref}{{[originallemma][0][]0F9R}{[1][0][]source}{}{}{}}
\makeatother
\crefname{originallemma}{original lemma}{original lemmas}
\Crefname{originallemma}{Original lemma}{Original lemmas}
\makeatletter
\newlabel{lemma-glueing-f}{{0F9S}{source}{}{sourcepointer-4}{}}
\newlabel{lemma-glueing-f@cref}{{[originallemma][0][]0F9S}{[1][0][]source}{}{}{}}
\makeatother
\crefname{originallemma}{original lemma}{original lemmas}
\Crefname{originallemma}{Original lemma}{Original lemmas}
\makeatletter
\newlabel{spaces-properties-lemma-points-cartesian}{{03H4}{source}{}{sourcepointer-5}{}}
\newlabel{spaces-properties-lemma-points-cartesian@cref}{{[originallemma][0][]03H4}{[1][0][]source}{}{}{}}
\makeatother
\title{The glueing functor is fully faithful on glueable spaces with affine diagonal}
\author{The Stacks Project Authors}
\date{stacks-project:spaces-pushouts.tex:local-snapshot-20261012; GNU FDL 1.2 or later}
\begin{document}
\maketitle
\noindent\textit{Editorial note.} Statement span is lines 2685-2692 of sources/stacks-spaces-pushouts/pinned-original-source.tex; proof span is lines 2694-2755, the complete author proof (60 body lines).\par\medskip
\noindent\textit{Editorial note.} Navigation stubs: 5 cited-but-not-included labels are supplied as non-load-bearing pointers with their pinned tags where the snapshot tag index has them.\par\medskip
\noindent\textit{Editorial note.} extra\_packages is empty on purpose: the retained preamble already loads xy with options and a second declaration clashes in the assembler wrapper.\par\medskip
\noindent\textit{Editorial note.} Text is retained verbatim, including the author's own typos and placeholders.\par\medskip
\noindent\textit{Editorial note.} Queue index 9 of queue\_spaces-pushouts.json (pb=2657, stmt 2685-2692, proof 2694-2755).\par\medskip
\noindent\textit{Editorial note.} spaces-pushouts wave-1: fresh adjudication accept (2026-10-14); six-check dispatched pool (§2 row appended 02:4x).\par\medskip
\noindent\textit{Allowed standard prerequisites (editorial).}\begin{itemize}
\item Explicit prerequisite boundary. The cited inputs are prerequisites; the argument in the retained span is not replaced by any cited result.
\end{itemize}\medskip
\ifdefined\hypertarget\hypertarget{sourcepointer-1}{}\fi
\noindent\textit{Source reference.} Original lemma 0F9N: sources/stacks-spaces-pushouts/pinned-original-source.tex, source lines 2470--2474. Cited by the retained statement or proof but not part of the collected spans; supplied only as a navigation pointer. The argument does not depend on its content.\par\medskip
\ifdefined\hypertarget\hypertarget{sourcepointer-2}{}\fi
\noindent\textit{Source reference.} Original lemma 0F9Q: sources/stacks-spaces-pushouts/pinned-original-source.tex, source lines 2547--2555. Cited by the retained statement or proof but not part of the collected spans; supplied only as a navigation pointer. The argument does not depend on its content.\par\medskip
\ifdefined\hypertarget\hypertarget{sourcepointer-3}{}\fi
\noindent\textit{Source reference.} Original lemma 0F9R: sources/stacks-spaces-pushouts/pinned-original-source.tex, source lines 2582--2589. Cited by the retained statement or proof but not part of the collected spans; supplied only as a navigation pointer. The argument does not depend on its content.\par\medskip
\ifdefined\hypertarget\hypertarget{sourcepointer-4}{}\fi
\noindent\textit{Source reference.} Original lemma 0F9S: sources/stacks-spaces-pushouts/pinned-original-source.tex, source lines 2648--2653. Cited by the retained statement or proof but not part of the collected spans; supplied only as a navigation pointer. The argument does not depend on its content.\par\medskip
\ifdefined\hypertarget\hypertarget{sourcepointer-5}{}\fi
\noindent\textit{Source reference.} Original lemma 03H4: sources/stacks-spaces-pushouts/pinned-original-source.tex, source lines 2755--2755. Cited by the retained statement or proof but not part of the collected spans; supplied only as a navigation pointer. The argument does not depend on its content.\par\medskip
\begin{lemma}
\label{lemma-glueing-ff}
Let $(R \to R', f)$ be a glueing pair, see above.
The functor (\ref{equation-beauville-laszlo-glueing-spaces})
is fully faithful on the full subcategory of
algebraic spaces $Y/X$ which are (a) glueable for $(R \to R', f)$
and (b) have affine diagonal $Y \to Y \times_X Y$.
\end{lemma}

\begin{proof}
Let $Y, Z$ be two algebraic spaces over $X$ which are both glueable for
$(R \to R', f)$ and assume the diagonal of $Z$ is affine. Let
$a : U \times_X Y \to U \times_X Z$ over $U$ and
$b : X' \times_X Y \to X' \times_X Z$ over $X'$ be two morphisms
of algebraic spaces
which induce the same morphism $c : U' \times_X Y \to U' \times_X Z$
over $U'$.
We want to construct a morphism $f : Y \to Z$ over $X$
which produces the morphisms $a$, $b$ on base change to $U$, $X'$.
By the faithfulness of Lemma \ref{lemma-glueing-f}, it suffices to construct
the morphism $f$ \'etale locally on $Y$ (details omitted).
Thus we may and do assume $Y$ is affine.

\medskip\noindent
Let $y \in |Y|$ be a point. If $y$ maps into the open $U \subset X$,
then $U \times_X Y$ is an open of $Y$ on which the morphism $f$
is defined (we can just take $a$).
Thus we may assume $y$ maps into the closed subset
$V(f)$ of $X$. Since $R/fR = R'/fR'$ there is a unique point
$y' \in |X' \times_X Y|$ mapping to $y$. Denote
$z' = b(y') \in |X' \times_X Z|$ and $z \in |Z|$ the images of $y'$.
Choose an \'etale neighbourhood $(W, w) \to (Z, z)$
with $W$ affine. Observe that
$$
(U \times_X W) \times_{U \times_X Z, a} (U \times_X Y),\quad
(U' \times_X W) \times_{U' \times_X Z, c} (U' \times_X Y),
$$
and
$$
(X' \times_X W) \times_{X' \times_X Z, b} (X' \times_X Y)
$$
form an object of $\textit{Spaces}(U \leftarrow U' \to X')$
with affine parts (this is where we use that $Z$ has affine diagonal).
Hence by Lemma \ref{lemma-glueing-affines}
there exists a unique affine scheme $V$ glueable for $(R \to R', f)$ such that
$$
(U \times_X V, U' \times_X V, X' \times_X V)
$$
is the triple displayed above. By fully faithfulness for the affine
case (Lemma \ref{lemma-glueing-affines}) we get a unique morphisms
$V \to W$ and $V \to Y$ agreeing with the first and second projection
morphisms over $U$ and $X'$ in the construction above.
By Lemma \ref{lemma-glueing-affines-etale} the morphism $V \to Y$ is \'etale.
To finish the proof, it suffices to show that there is a point $v \in |V|$
mapping to $y$ (because then $f$ is defined on an \'etale neighbourhood
of $y$, namely $V$).
There is a unique point $w' \in |X' \times_X W|$ mapping to $w$.
By uniqueness $w'$ is mapped to $z'$ under the map
$|X' \times_X W| \to |X' \times_X Z|$. Then we consider the cartesian
diagram
$$
\xymatrix{
X' \times_X V \ar[r] \ar[d] & X' \times_X W \ar[d] \\
X' \times_X Y \ar[r] & X' \times_X Z
}
$$
to see that there is a point $v' \in |X' \times_X V|$
mapping to $y'$ and $w'$, see
Properties of Spaces, Lemma \ref{spaces-properties-lemma-points-cartesian}.
Of course the image $v$ of $v'$ in $|V|$ maps to $y$ and the proof is complete.
\end{proof}

\end{document}
